\begin{table}[t]
\centering
\caption{公开数据与 Nishimura 等 \cite{nishimura2026} 报告统计量的一致性核对，核对量涉及观测数、残差平方和及案例组成。}\label{tab:datacheck}
\scriptsize
\begin{adjustbox}{max width=\linewidth}
\begin{tabular}{p{9.2cm}rrl}
\toprule
核对项 & 文献报告 & 公开数据 & 是否一致 \\
\midrule
文献表 1 Akerman70 全模型，用公开数据的 nT=36 复算文献报告 BIC（Fobj=59.0，k=5） & 35.7 & 35.7 & 是 \\
文献表 1 Akerman100 全模型，用公开数据的 nT=28 复算文献报告 BIC（Fobj=117.9，k=5） & 56.9 & 56.9 & 是 \\
文献表 2 Akerman70 简化模型，用公开数据的 nT=36 复算文献报告 BIC（Fobj=70.9，k=3） & 35.1 & 35.1 & 是 \\
文献表 2 Akerman100 简化模型，用公开数据的 nT=28 复算文献报告 BIC（Fobj=118.0，k=3） & 50.3 & 50.3 & 是 \\
文献表 3 Bornadel 全模型，用公开数据的 nT=120 复算文献报告 BIC（Fobj=329.9，k=5） & 145.3 & 145.3 & 是 \\
文献表 3 Bornadel 简化模型，用公开数据的 nT=120 复算文献报告 BIC（Fobj=329.9，k=4） & 140.5 & 140.5 & 是 \\
文献表 4 Tao025 全模型，用公开数据的 nT=64 复算文献报告 BIC（Fobj=121.2，k=5） & 61.7 & 61.7 & 是 \\
文献表 4 Tao035+045 全模型，用公开数据的 nT=144 复算文献报告 BIC（Fobj=325.8，k=5） & 142.4 & 142.4 & 是 \\
文献表 5 Tao025 简化模型，用公开数据的 nT=64 复算文献报告 BIC（Fobj=123.5，k=3） & 54.5 & 54.5 & 是 \\
文献表 5 Tao035+045 简化模型，用公开数据的 nT=144 复算文献报告 BIC（Fobj=326.7，k=3） & 132.9 & 132.9 & 是 \\
文献表 6 Mao1-3 全模型，用公开数据的 nT=141 复算文献报告 BIC（Fobj=1475.9，k=7） & 365.7 & 365.7 & 是 \\
文献表 6 Mao1-3 简化模型，用公开数据的 nT=141 复算文献报告 BIC（Fobj=1771.5，k=5） & 381.6 & 381.6 & 是 \\
文献来源 / 案例数 & 4 & 4 & 是 \\
Åkerman 数据集（70 和 100 °C） & 2 & 2 & 是 \\
Tao 数据集（aw 0.25/0.35/0.45） & 3 & 3 & 是 \\
Mao 批次（列出 / 使用） & 5 / 3 & 5 / 3 & 是 \\
指纹拟合 Akerman，用公开数据核对残差平方和 & 74.2 & 74.2 & 是 \\
指纹拟合 Bornadel，用公开数据核对残差平方和 & 840.6 & 840.6 & 是 \\
指纹拟合 Tao，用公开数据核对残差平方和 & 535.3 & 535.3 & 是 \\
指纹拟合 Mao，用公开数据核对残差平方和 & 1616.0 & 1616.0 & 是 \\
\bottomrule
\end{tabular}
\end{adjustbox}
\end{table}
